This book is intended to introduce its readers to the mathematical theory of partial differential equations. But to suggest that there is a "theory" of partial differential equations (in the same sense that there is a theory of ordinary differential equations or a theory of functions of a single complex variable) is misleading. PDEs is a much larger subject than the two mentioned above (it includes both of them as special cases) and a less well developed one. However, although a casual observer may decide the subject is simply a grab bag of unrelated techniques used to handle different types of problems, there are in fact certain themes that run throughout.

In order to illustrate these themes we take two approaches. The first is to pose a group of questions that arise in many problems in PDEs (existence, multiplicity, etc.). As examples of different methods of attacking these problems, we examine some results from the theories of ODEs, advanced calculus and complex variables (with which the reader is assumed to have some familiarity). The second approach is to examine three partial differential equations (Laplace's equation, the heat equation and the wave equation) in a very elementary fashion (again, this will probably be a review for most readers). We will see that even the most elementary methods foreshadow deeper results found in the later chapters of this book.

\paragraph*{Nonuniqueness for linear and nonlinear boundary-value problems}

While uniqueness is often a desirable property for a solution of a problem (often for physical reasons), there are situations in which multiple solutions are desirable. A common mathematical problem involving multiple solutions is an eigenvalue problem. The reader should, of course, be familiar with the various existence and multiplicity results from finite-dimensional linear algebra, but let us consider a few problems from ordinary differential equations. We consider the following second-order ODE depending on the parameter $\lambda$:
\begin{equation}\label{eq:linear-mixed-boundary-ode}
u''+\lambda u=0.
\end{equation}

Of course, if we imposed two initial conditions (at one point in space) Theorem~\ref{thm:ode-uniqueness} would imply that we would have a unique solution. (To apply the theorem directly we need to convert the problem from a second-order equation to a first-order system.) However, if we impose the two-point boundary conditions
\begin{equation}\label{eq:linear-mixed-dirichlet}
u(0)=0,
\end{equation}
\begin{equation}\label{eq:linear-mixed-neumann}
u'(1)=0,
\end{equation}
the uniqueness theorem does not apply. Instead we get the following result.

\begin{Theorem}
\label{thm:linear-mixed-boundary-alternatives}
There are two alternatives for the solutions of the boundary-value problem \eqref{eq:linear-mixed-boundary-ode}, \eqref{eq:linear-mixed-dirichlet}, \eqref{eq:linear-mixed-neumann}.
\begin{enumerate}
\item If $\lambda=\lambda_n:=((2n+1)^2\pi^2)/4$, $n=0,1,2,\ldots$, then the boundary-value problem has a family of solutions parameterized by $A\in(-\infty,\infty)$:
\[
u_n(x)=A\sin\dfrac{(2n+1)\pi}{2}x.
\]
In this case we say $\lambda$ is an eigenvalue.
\item For all other values of $\lambda$ the only solution of the boundary-value problem is the trivial solution
\[
u(x)\equiv0.
\]
\end{enumerate}
\end{Theorem}

This characteristic of having either a unique (trivial) solution or an infinite linear family of solutions is typical of linear problems. More interesting multiplicity results are available for nonlinear problems and are the main subject of modern \emph{bifurcation theory}. For example, consider the following nonlinear boundary-value problem, which was derived by Euler to describe the deflection of a thin, uniform, inextensible, vertical, elastic beam under a load $\lambda$:
\begin{equation}\label{eq:euler-elastica}
\theta''(x)+\lambda\sin\theta(x)=0,
\end{equation}
\begin{equation}\label{eq:euler-elastica-dirichlet}
\theta(0)=0,
\end{equation}
\begin{equation}\label{eq:euler-elastica-neumann}
\theta'(1)=0.
\end{equation}

\begin{figure}[H]
\centering
\includegraphics[width=.78\linewidth]{figures/beam-bifurcation-scan.png}
\caption{Bifurcation diagram for the nonlinear boundary-value problem}\label{fig:beam-bifurcation}
\end{figure}

(Note that the linear ODE \eqref{eq:linear-mixed-boundary-ode} is an approximation of \eqref{eq:euler-elastica} for small $\theta$.)

Solutions of this nonlinear boundary-value problem have been computed in closed form (in terms of Jacobi elliptic functions) and are probably best displayed by a \emph{bifurcation diagram} such as Figure~\ref{fig:beam-bifurcation}. This figure displays the amplitude of a solution $\theta$ as a function of the value of $\lambda$ at which the solution occurs. The $\lambda$ axis denotes the trivial solution $\theta\equiv0$ (which holds for every $\lambda$). Note that a \emph{branch} of nontrivial solutions emanates from each of the eigenvalues of the linear problem above. Thus for $\lambda\in(\lambda_{n-1},\lambda_n)$, $n=1,2,3,\ldots$, there are precisely $2n$ nontrivial solutions of the boundary-value problem.

\subsubsection{Stability}\label{subsec:stability}

The term stability is one that has a variety of different meanings within mathematics. One often says that a problem is stable if it is "continuous with respect to the data"; i.e., a problem is stable if when we change the problem "slightly," the solution changes only slightly. We make this precise below in the context of initial-value problems for ODEs. Another notion of stability is that of "asymptotic stability." Here we say a problem is stable if all of its solutions get close to some "nice" solution as time goes to infinity. We make this notion precise with a result on linear systems of ODEs with constant coefficients.

\paragraph*{Stability with respect to initial conditions}

In this section we assume that $\bm{F}$ satisfies the hypotheses of Theorem~\ref{thm:ode-existence}, and we define $\hat{\bm{y}}(t,t_0,\bm{y}_0)$ to be the unique solution of \eqref{eq:ode-initial-value-dynamics}, \eqref{eq:ode-initial-value-data}. We then have the following standard result.

\begin{Theorem}{Continuity with respect to initial conditions}
\label{thm:initial-data-continuity}
The function $\hat{\bm{y}}$ is well defined on an open set
\[
U\subset\mathbb{R}\times D.
\]
Furthermore, at every $(t,t_0,\bm{y}_0)\in U$ the function
\[
(t_0,\bm{y}_0)\mapsto\hat{\bm{y}}(t,t_0,\bm{y}_0)
\]
is continuous; i.e., for any $\epsilon>0$ there exists $\delta$ (depending on $(t,t_0,\bm{y}_0)$ and $\epsilon$) such that if
\[
|(t_0,\bm{y}_0)-(\tilde{t}_0,\tilde{\bm{y}}_0)|<\delta,
\]
then $\hat{\bm{y}}(t,\tilde{t}_0,\tilde{\bm{y}}_0)$ is well defined and
\begin{equation}\label{eq:initial-data-continuity}
|\hat{\bm{y}}(t,t_0,\bm{y}_0)-\hat{\bm{y}}(t,\tilde{t}_0,\tilde{\bm{y}}_0)|<\epsilon.
\end{equation}
\end{Theorem}

Thus, we see that small changes in the initial conditions result in small changes in the solutions of the initial-value problem.
